% Résumé of Saif Ur Rehman. Build with `npm run resume` (resume/build.sh).
% Keep every claim in line with src/data/: statuses, numbers and roles must
% match what the site says.
\documentclass[10pt]{article}
% Layout for saif-ur-rehman-resume.tex. Black text on white, Charter,
% small-caps section rules, bold titles with right-aligned dates.
\usepackage[a4paper,left=1.35cm,right=1.35cm,top=1.1cm,bottom=1.1cm]{geometry}
\usepackage[T1]{fontenc}
\usepackage[utf8]{inputenc}
\usepackage{XCharter}
\usepackage[protrusion=true,expansion=true]{microtype}
\usepackage{titlesec}
\usepackage{enumitem}
\usepackage{xcolor}
\usepackage{hyperref}
\hypersetup{
  hidelinks,
  pdftitle={Saif Ur Rehman - Resume},
  pdfauthor={Saif Ur Rehman},
  pdfsubject={Machine learning engineer and AI researcher},
  pdflang={en-GB},
  pdfdisplaydoctitle=true,
}

\pagestyle{empty}
\setlength{\parindent}{0pt}
\setlength{\parskip}{0pt}
\linespread{1.0}
\color{black}

\titleformat{\section}{\large\scshape}{}{0pt}{}[\vspace{-0.9ex}\rule{\textwidth}{0.4pt}]
\titlespacing*{\section}{0pt}{1.25ex}{0.8ex}

\setlist[itemize]{leftmargin=1.05em,label={\small\textbullet},itemsep=0.1ex,topsep=0.3ex,parsep=0pt,partopsep=0pt}

% \entry{title}{dates}{left meta}{right meta}
\newcommand{\entry}[4]{%
  \par\textbf{#1}\hfill{\small #2}\par\nopagebreak
  \if\relax\detokenize{#3#4}\relax\else{\small\itshape #3}\hfill{\small\itshape #4}\par\nopagebreak\fi}
% \project{title}{link text}{url}{stack}; with an empty url the link text is plain.
\newcommand{\project}[4]{%
  \par\textbf{#1}\hfill{\small\if\relax\detokenize{#3}\relax #2\else\href{#3}{#2}\fi}\par\nopagebreak
  \if\relax\detokenize{#4}\relax\else{\small\itshape #4}\par\nopagebreak\fi}
\newcommand{\gap}{\par\vspace{0.55ex}}
% \skill{label}{items}: one line of the skills list, wrapped lines indented.
\newcommand{\skill}[2]{\par\hangindent=1em\hangafter=1\textbf{#1:} #2\par}


\begin{document}

{\centering
{\LARGE\scshape Saif Ur Rehman}\par\vspace{0.6ex}
{\small Liverpool, UK \enspace\textbullet\enspace
\href{mailto:saif10urrehman@gmail.com}{saif10urrehman@gmail.com} \enspace\textbullet\enspace
\href{https://www.linkedin.com/in/saifurrhmn/}{linkedin.com/in/saifurrhmn} \enspace\textbullet\enspace
\href{https://github.com/saifrhman}{github.com/saifrhman} \enspace\textbullet\enspace
\href{https://saifrhman.github.io}{saifrhman.github.io}}\par}

\vspace{1.2ex}
MSc Data Science and AI candidate, University of Liverpool. Research on geometry-aware curation of protein training
data (dissertation), evidence use in feed-forward 4D reconstruction (first-author manuscript in development) and
peer-reward misalignment in multi-agent LLMs (workshop submission, under review). Builds and evaluates ML systems in
Python and PyTorch; applied AI intern at Rofsix, after systems, cloud and support engineering roles (2020--2025).

\section{Education}
\entry{MSc Data Science and Artificial Intelligence, University of Liverpool}{Sep 2025 -- 2026}{}{}
Dissertation: \emph{Beyond AlphaFold: Filling the Blind Spots in Protein Structure Prediction}, supervised by
Dr Olga Anosova. Modules include Machine Learning and Bio-inspired Optimisation, Computational Intelligence and
Applied AI.
\gap
\entry{BE Electrical (Telecommunication) Engineering, NUST, Pakistan}{2016 -- 2021}{}{}

\section{Research experience}
\entry{Beyond AlphaFold / PDBClean: Geometry-Aware Curation of Protein Structure Training Data}{2026}
  {MSc dissertation, sole author \textbullet\ Python, PyTorch, OpenFold, MMseqs2, gemmi, Slurm}
  {\href{https://github.com/saifrhman/702_BeyondAF}{github.com/saifrhman/702\_BeyondAF}}
\begin{itemize}
  \item Re-implemented and scaled the published Backbone Rigid Invariant (BRI) curation protocol as PDBClean, a
    configuration-driven pipeline over a fixed PDB snapshot (1 Jan 2026). A full re-run reproduced its output exactly.
  \item Implemented Elkin and Kurlin's (2023) compressed cover tree for exact $L_\infty$ search, behind a length
    partition and a lossless lower-bound prefilter.
  \item Removed 78,754 geometric near-duplicates ($\le 0.010$\,\AA) from 578,524 eligible chains, keeping 499,770.
    The cover tree's output matched an exhaustive check of all 3.53M candidate pairs, pair for pair.
  \item Found the geometric and sequence criteria nearly nested: 99.5\% of geometric removals are sequence-identical
    to their representative (MMseqs2 on the same population).
  \item Prepared OpenFold inputs with fresh MSAs for a 142,056-chain, one-per-sequence set. From-scratch retraining
    under Slurm is in progress, with no accuracy results yet; the training scripts were adapted from a predecessor project.
\end{itemize}
\gap
\entry{Separating Solvers and Evaluators Is Not Enough: Peer-Reward Misalignment in Multi-Agent LLMs}{2026}
  {Co-author, four-person team \textbullet\ submitted to a NeurIPS 2026 workshop, under review}{}
\begin{itemize}
  \item Two LoRA policies on a shared frozen 7--9B base each solve a task and grade the other's answer. In most
    matched pairs, rewarding the peer's grade instead of the task metric lowered external task performance.
  \item Wrote the single-agent RL training loop (REINFORCE with a moving-average baseline) that the two-agent trainer
    builds on, and produced the single-agent replication figure.
  \item Follow-up: audited how policy-gradient credit reaches the grading tokens; designed a reward-topology
    comparison (objective reward, frozen judge, four-agent cycle, reciprocal pairs), not yet run.
\end{itemize}
\gap
\entry{Geometric Evidence Utilization in Feed-Forward 4D Reconstruction (working title)}{2026 -- present}
  {First-author manuscript in development \textbullet\ research in progress \textbullet\ Target: CVPR 2027}{}
\begin{itemize}
  \item Asks whether feed-forward 4D reconstruction models use stronger geometric conditioning to estimate a
    physical quantity correctly, or fall back on a learned prior where the input does not determine it.
  \item Built a synthetic benchmark with an exact ambiguity endpoint (physically different worlds render
    byte-identical video) and a classical reference estimator. A ground-truth camera-pose intervention on Any4D and
    4RC is in progress; no findings yet.
  \item Designing KINETIC, which represents dynamic scenes as independently transforming kinematic components with a
    separate global gauge. One transformation family and 20 test scenes are held out from all development.
\end{itemize}
\gap
\entry{Geographical Data and Computational Tools for Determining Restoration Priorities in Mangrove Ecosystems: A Scoping Review}{2025 -- 2026}
  {Co-author, five-person team \textbullet\ manuscript in preparation \textbullet\ began as MSc group coursework}{}
PRISMA-ScR review of 24 included studies. Restructured the manuscript and drafted large parts of the data-source,
sensing-modality, methods and future-work sections.

\section{Experience}
\entry{Applied AI Intern, Rofsix (UK)}{2026 -- present}{}{}
\begin{itemize}
  \item Developing a real-time system that detects and redacts confidential information during screen sharing,
    combining OCR, transformer-based NER and image processing. Built and evaluated inference components with ONNX,
    OpenVINO and PaddleOCR, and profiled the pipeline for low-latency CPU execution.
\end{itemize}
\gap
\entry{Senior Systems Engineer, RoshTech}{Nov 2022 -- Sep 2025}{}{}
\begin{itemize}
  \item Built threat-detection and anomaly-analysis workflows over MongoDB, Graylog and the Elastic Stack to support
    incident investigation for enterprise clients.
\end{itemize}
\gap
\entry{Data Science Intern (part-time), Creative Tech Solutions (Islamabad)}{Jan 2024 -- Apr 2024}{}{}
\begin{itemize}
  \item Wrote Python scripts for data analysis and data cleaning for the analytics team.
\end{itemize}
\gap
\entry{Earlier roles}{2020 -- 2022}
  {Senior Cloud Engineer, Tier3tech (Feb 2022 -- Nov 2022) \textbullet\ Level 2 Support Team Lead, MiGo Innovations (Nov 2020 -- Feb 2022)}{}

\section{Selected projects}
\project{F1 AI Copilot: Retrieval-Augmented QA over the FIA Regulations}{github.com/saifrhman/f1-ai-copilot}
  {https://github.com/saifrhman/f1-ai-copilot}{Python \textbullet\ LangChain \textbullet\ Qdrant \textbullet\ FastAPI \textbullet\ Streamlit \textbullet\ pytest}
\begin{itemize}
  \item Question answering over the 2026 FIA F1 regulations (592 pages). A deterministic validator declines any
    answer whose citations, rule identifiers or figures are not in the passages it cites (sentences marked as
    inferences are exempt from the number check). The index fingerprints its build inputs; 1,442 tests run in CI.
\end{itemize}
\gap
\project{Chest X-ray Classification (group coursework)}{COMP534 Applied AI, University of Liverpool}{}
  {PyTorch \textbullet\ EfficientNet-B3 \textbullet\ Swin V2-S \textbullet\ MaxViT-T \textbullet\ stratified 5-fold CV}
\begin{itemize}
  \item Team of four; 3,475 course-provided chest X-rays (Normal, Opacity, viral Pneumonia). Compared a custom
    116,931-parameter CNN with EfficientNet-B3, Swin V2-S and MaxViT-T, each from scratch, frozen and fine-tuned.
  \item From scratch Swin V2-S failed to train (CV macro-F1 0.29); fine-tuned it scored highest (0.953) and reached
    test macro-F1 0.948, Cohen's $\kappa$ 0.92 on 695 held-out images. Single run, not clinically validated; the
    Pneumonia class likely differs from the others in patient age and image source.
\end{itemize}
\gap
\project{Hybrid Receipt Extraction: OCR, Rules and LayoutLMv3}{github.com/saifrhman/noise-robust-ocr-pipeline}
  {https://github.com/saifrhman/noise-robust-ocr-pipeline}{EasyOCR \textbullet\ LayoutLMv3 \textbullet\ PyTorch \textbullet\ Transformers \textbullet\ weak supervision}
\begin{itemize}
  \item EasyOCR, a rule parser and a LayoutLMv3 token classifier, merged field by field. SROIE labels only 4 fields,
    so rule-derived pseudo-labels for 13 more entity types are down-weighted (0.76--0.86 against 1.0) in the token
    loss. No evaluation results have been committed yet.
\end{itemize}
\gap
\project{Football Intelligence Platform}{github.com/saifrhman/football-analytics-platform}
  {https://github.com/saifrhman/football-analytics-platform}{Python \textbullet\ BigQuery \textbullet\ dbt \textbullet\ FastAPI \textbullet\ Streamlit \textbullet\ GitHub Actions}
\begin{itemize}
  \item Pipeline from StatsBomb Open Data to BigQuery: bronze raw JSON, nine typed silver tables with explicit
    schemas, and dbt gold models (5 dimensions, 4 facts, 23 schema tests), served through FastAPI and Streamlit. Run
    end to end on a five-match sample; CI runs 39 unit tests, ruff and dbt parse.
\end{itemize}
\gap
\project{TacticLens: Structural Retrieval over Football Tracking Data}{github.com/saifrhman/tacticlens}
  {https://github.com/saifrhman/tacticlens}{Research design and foundation code \textbullet\ Python \textbullet\ Pydantic \textbullet\ FastAPI}
\begin{itemize}
  \item Research design (questions, baselines and rejection criteria fixed in advance) for retrieving passages of
    play with similar multi-player movement, and a tracking data model with seven typed observation statuses. The
    retrieval model, index and feedback loop are not yet built.
\end{itemize}
\gap
\project{Sample-Efficient Learning: Active Learning and Bayesian Optimisation}{}{}
  {scikit-learn \textbullet\ BoTorch \textbullet\ GPyTorch \textbullet\
  \href{https://github.com/saifrhman/active-learning-classification}{active-learning-classification} \textbullet\
  \href{https://github.com/saifrhman/closed-loop-hpo-botorch}{closed-loop-hpo-botorch}}
\begin{itemize}
  \item Entropy-based active learning reached 0.986 test accuracy after 60 labels, against 70 for random acquisition.
    GP Bayesian optimisation (Expected Improvement) reached best 5-fold CV accuracy 0.961 against 0.956 for random
    search in 33 random-forest tuning evaluations. Single seed each.
\end{itemize}

\section{Competitive programming and hackathons}
\skill{ICPC 2025}{University of Liverpool team ``Baby Harp Seal'': NWERC 2025, 107th; ICPC UK \& Ireland
  Programming Contest 2025, 50th overall and 1st at the University of Liverpool site.}
\skill{Hackathons}{Liverpool Computer Society Hackathon, 1st place (team), deep-learning solution; Ultamation
  Hackathon, 2nd place, constraint-satisfaction scheduling.}

\section{Skills and certifications}
\skill{Languages}{Python, C++, SQL, Bash, R}
\skill{LLMs and RL fine-tuning}{LoRA/QLoRA, REINFORCE, multi-agent LLM training with peer-graded rewards,
  LLM-as-judge rewards, retrieval-augmented generation (LangChain, Qdrant), Hugging Face Transformers}
\skill{ML training and evaluation}{PyTorch, transfer learning (EfficientNet, Swin V2, MaxViT), Slurm GPU training
  and checkpointing (OpenFold), large-scale dataset curation and deduplication, cross-validation, error analysis,
  active learning, Bayesian optimisation (BoTorch, GPyTorch)}
\skill{Geometry, vision and science}{3D/4D reconstruction models, protein structure data (OpenFold, MMseqs2,
  gemmi, Mol*), exact metric search (cover trees, k-d trees), LayoutLMv3, EasyOCR, PaddleOCR, OpenCV, YOLOv8, U-Net}
\skill{Systems and data}{HPC clusters (Slurm), OpenMP, ONNX, OpenVINO, FastAPI, Streamlit, BigQuery, dbt,
  PostgreSQL, MongoDB, Elastic Stack, Azure, Docker, Git, GitHub Actions, Linux}

\skill{Certifications}{Fundamentals of Accelerated Computing with Modern CUDA C++, NVIDIA (2026); Machine Learning
  Specialization, DeepLearning.AI on Coursera (2025); IBM Data Science Professional Certificate, IBM on Coursera
  (2025); Associate Data Engineer in SQL track, DataCamp (2025)}

\end{document}
